
\documentclass{article}
\usepackage{colortbl}
\usepackage{makecell}
\usepackage{multirow}
\usepackage{supertabular}

\begin{document}

\newcounter{utterance}

\centering \large Interaction Transcript for game `cladder', experiment `full\_v1.5\_default', episode 6482 with refineonly{-}Qwen3.5{-}9B.
\vspace{24pt}

{ \footnotesize  \setcounter{utterance}{1}
\setlength{\tabcolsep}{0pt}
\begin{supertabular}{c@{$\;$}|p{.15\linewidth}@{}p{.15\linewidth}p{.15\linewidth}p{.15\linewidth}p{.15\linewidth}p{.15\linewidth}}
    \# & \multicolumn{2}{c}{Player} && \multicolumn{2}{c}{Game Master} \\
    \hline

    \theutterance \stepcounter{utterance}  
    & & & \multicolumn{4}{p{0.6\linewidth}}{
        \cellcolor[rgb]{0.9,0.9,0.9}{
            \makecell[{{p{\linewidth}}}]{
                \texttt{\tiny{[P1$\langle$GM]}}
                \texttt{You are an expert in causal inference. The following question is not a typical commonsense query, but rather a meticulously designed question created by a professor specializing in causal inference, intended to assess the students' mastery of the course content.} \\
\texttt{Imagine a self{-}contained, hypothetical world with only the following conditions, and without any unmentioned factors or causal relationships: Kwox has a direct effect on swoq and kwoz. Swoq has a direct effect on kwoz. For those who are not kwox, the probability of kwoz is 36\%. For those who are kwox, the probability of kwoz is 56\%. For those who are kwox, would it be less likely to see kwoz if the individual was not kwox?} \\
\texttt{Start your answer with 'Yes' or 'No', followed by additional reasoning or evidence to support your explanation.} \\
            }
        }
    }
     \\ \\

    \theutterance \stepcounter{utterance}  
    & & \multicolumn{4}{p{0.6\linewidth}}{
        \cellcolor[rgb]{0.95,0.95,0.95}{
            \makecell[{{p{\linewidth}}}]{
                \texttt{\tiny{[P1$\leftrightarrow$P]}}
                \texttt{{[}DRAFT{]}} \\
\texttt{No} \\
\\ 
\texttt{**Reasoning:**} \\
\\ 
\texttt{To answer this question, we must carefully parse the logical structure of the counterfactual being asked. The question asks: "For those who are kwox, would it be less likely to see kwoz if the individual was not kwox?"} \\
\\ 
\texttt{This is a counterfactual query regarding the effect of the variable \$Kwox\$ on the outcome \$Kwoz\$. Let us define the probabilities based on the provided conditions:} \\
\texttt{1.  **Factual State (\$Kwox = 1\$):** The probability of \$Kwoz\$ is given as \$P(Kwoz | Kwox) = 56\textbackslash{}\%\$.} \\
\texttt{2.  **Counterfactual State (\$Kwox = 0\$):** The question asks for the probability of \$Kwoz\$ for the same individuals (who are currently \$Kwox\$) had they *not* been \$Kwox\$. This is denoted as \$P(Kwoz | do(Kwox=0))\$ or, under the assumption of no unmentioned confounders in this self{-}contained world, \$P(Kwoz | \textbackslash{}neg Kwox)\$.} \\
\\ 
\texttt{The prompt explicitly states: "For those who are not kwox, the probability of kwoz is 36\%." Therefore, the counterfactual probability is \$36\textbackslash{}\%\$.} \\
\\ 
\texttt{The question asks if the likelihood of seeing \$Kwoz\$ would be **less** in the counterfactual scenario (\$36\textbackslash{}\%\$) compared to the factual scenario (\$56\textbackslash{}\%\$).} \\
\texttt{*   Factual Probability: \$56\textbackslash{}\%\$} \\
\texttt{*   Counterfactual Probability: \$36\textbackslash{}\%\$} \\
\\ 
\texttt{Since \$36\textbackslash{}\% < 56\textbackslash{}\%\$, the probability of observing \$Kwoz\$ is indeed lower if the individual was not \$Kwox\$.} \\
\\ 
\texttt{However, there is a potential linguistic trap in the phrasing "For those who are kwox... if the individual was not kwox."} \\
\texttt{If the question implies a comparison of the *marginal* probabilities across the population defined by the condition "For those who are kwox" (which is confusing phrasing), we must look at the causal effect.} \\
\texttt{The causal effect of \$Kwox\$ on \$Kwoz\$ is \$P(Kwoz | Kwox) {-} P(Kwoz | \textbackslash{}neg Kwox) = 56\textbackslash{}\% {-} 36\textbackslash{}\% = +20\textbackslash{}\%\$.} \\
\texttt{Because the effect is positive, removing the cause (\$Kwox\$) reduces the probability of the effect (\$Kwoz\$).} \\
\\ 
\texttt{Therefore, for an individual who is \$Kwox\$, the probability of \$Kwoz\$ would be lower (\$36\textbackslash{}\%\$) in the world where they are not \$Kwo} \\
            }
        }
    }
    & \\ \\

    \theutterance \stepcounter{utterance}  
    & \multicolumn{4}{p{0.6\linewidth}}{
        \cellcolor[rgb]{0.561,0.737,0.561}{
            \makecell[{{p{\linewidth}}}]{
                \texttt{\tiny{[P1$\rangle$GM]}}
                \texttt{No} \\
\\ 
\texttt{**Reasoning:**} \\
\\ 
\texttt{To answer this question, we must carefully parse the logical structure of the counterfactual being asked. The question asks: "For those who are kwox, would it be less likely to see kwoz if the individual was not kwox?"} \\
\\ 
\texttt{This is a counterfactual query regarding the effect of the variable \$Kwox\$ on the outcome \$Kwoz\$. Let us define the probabilities based on the provided conditions:} \\
\texttt{1.  **Factual State (\$Kwox = 1\$):** The probability of \$Kwoz\$ is given as \$P(Kwoz | Kwox) = 56\textbackslash{}\%\$.} \\
\texttt{2.  **Counterfactual State (\$Kwox = 0\$):** The question asks for the probability of \$Kwoz\$ for the same individuals (who are currently \$Kwox\$) had they *not* been \$Kwox\$. This is denoted as \$P(Kwoz | do(Kwox=0))\$ or, under the assumption of no unmentioned confounders in this self{-}contained world, \$P(Kwoz | \textbackslash{}neg Kwox)\$.} \\
\\ 
\texttt{The prompt explicitly states: "For those who are not kwox, the probability of kwoz is 36\%." Therefore, the counterfactual probability is \$36\textbackslash{}\%\$.} \\
\\ 
\texttt{The question asks if the likelihood of seeing \$Kwoz\$ would be **less** in the counterfactual scenario (\$36\textbackslash{}\%\$) compared to the factual scenario (\$56\textbackslash{}\%\$).} \\
\texttt{*   Factual Probability: \$56\textbackslash{}\%\$} \\
\texttt{*   Counterfactual Probability: \$36\textbackslash{}\%\$} \\
\\ 
\texttt{Since \$36\textbackslash{}\% < 56\textbackslash{}\%\$, the probability of observing \$Kwoz\$ is indeed lower if the individual was not \$Kwox\$.} \\
\\ 
\texttt{Therefore, for an individual who is \$Kwox\$, the probability of \$Kwoz\$ would be lower (\$36\textbackslash{}\%\$) in the world where they are not \$Kwox\$.} \\
            }
        }
    }
    & & \\ \\

    \theutterance \stepcounter{utterance}  
    & & & \multicolumn{2}{p{0.3\linewidth}}{
        \cellcolor[rgb]{0.9,0.9,0.9}{
            \makecell[{{p{\linewidth}}}]{
                \texttt{\tiny{[GM$|$GM]}}
                \texttt{No} \\
            }
        }
    }
    & & \\ \\

    \theutterance \stepcounter{utterance}  
    & & & \multicolumn{2}{p{0.3\linewidth}}{
        \cellcolor[rgb]{0.9,0.9,0.9}{
            \makecell[{{p{\linewidth}}}]{
                \texttt{\tiny{[GM$|$GM]}}
                \texttt{yes} \\
            }
        }
    }
    & & \\ \\

    \theutterance \stepcounter{utterance}  
    & & & \multicolumn{2}{p{0.3\linewidth}}{
        \cellcolor[rgb]{0.9,0.9,0.9}{
            \makecell[{{p{\linewidth}}}]{
                \texttt{\tiny{[GM$|$GM]}}
                \texttt{game\_result = LOSE} \\
            }
        }
    }
    & & \\ \\

\end{supertabular}
}

\end{document}
